\addcontentsline{toc}{section}{Introduction}
\section*{Introduction}

Ce chapitre décrit la démarche suivie, expose la méthode du futur audit du réseau du CNHU-HKM et les éléments d'architecture documentés par le service, compare les solutions IDS/IPS candidates, puis expose la conception du système de détection retenu. Les éléments de conception de la section~\ref{sec:conception} sont des propositions de travail, à valider avec le tuteur avant leur mise en œuvre.

\section{Méthodologie}

\subsection{Étapes du projet}

Le projet s'organise en cinq phases successives, étalées sur \acompleter{X} semaines : un audit du réseau existant, une phase de conception de l'architecture cible, une phase de réalisation en environnement de test virtualisé, une phase d'évaluation expérimentale par des scénarios d'attaque, et une phase de synthèse donnant lieu à des recommandations pour un déploiement progressif au CNHU-HKM. La phase de conception (phase 2) est réalisée et documentée dans ce mémoire. La phase d'audit (phase 1) doit être conduite selon la méthode décrite dans la suite de ce chapitre ; ses résultats ne sont pas présentés ici. Les phases 3 et 4 n'ont pas été exécutées : elles font l'objet, au chapitre~\ref{chap:3}, d'un protocole entièrement spécifié et prêt à mettre en œuvre. Le tableau~\ref{tab:calendrier} fixe le calendrier prévisionnel de ces phases ; les dates et durées effectives doivent être reportées depuis le journal de bord, aucune durée n'étant indiquée tant qu'elle n'est pas attestée.

\begin{table}[htbp]
\centering
\small
\begin{tabularx}{\textwidth}{cllX}
\hline
\textbf{Phase} & \textbf{Semaines effectives} & \textbf{Contenu} & \textbf{Livrable} \\
\hline
1 & \acompleter{dates} & Audit du réseau existant, découverte d'hôtes et de services & Rapport d'audit et constats \\
2 & \acompleter{dates} & Conception de l'architecture cible, zones et flux & Plan d'adressage, matrice de flux \\
3 & \acompleter{dates} & Réalisation du laboratoire virtualisé (OPNsense, Suricata, Wazuh) & Plateforme de test documentée \\
4 & \acompleter{dates} & Évaluation expérimentale : scénarios S1 à S7, états A, B et C & Tableau de résultats et preuves \\
5 & \acompleter{dates} & Synthèse, recommandations de déploiement progressif & Rapport final \\
\hline
\end{tabularx}
\caption{Calendrier prévisionnel de réalisation ; les semaines effectives sont à reporter depuis le journal de bord}
\label{tab:calendrier}
\end{table}

\subsection{Cadre éthique et autorisations}

Les manipulations du futur audit réseau (découverte d'hôtes et balayage de ports) devront être menées dans le respect d'un cadre formalisé avec le service informatique et réseau du CNHU-HKM. Le périmètre autorisé pour la consultation de la documentation réseau et pour les opérations de découverte sera fixé par une autorisation écrite du tuteur de stage \acompleter{joindre le document signé en annexe~H}, qui rappellera également l'interdiction de tout test d'intrusion actif sur le réseau de production sans accord préalable explicite. Ce cadre impose qu'aucune donnée patient ne soit consultée ni manipulée. Les informations sensibles relatives à l'infrastructure (adresses IP réelles, noms de machines) seront anonymisées lorsque leur divulgation n'apportera pas de valeur pédagogique. Les scénarios d'attaque du chapitre~\ref{chap:3} sont, quant à eux, exclusivement prévus dans le laboratoire virtualisé : ils ne doivent pas être exécutés sur le réseau de production.

\section{Audit du réseau : méthode envisagée et éléments documentés}

\subsection{Méthode de découverte envisagée}

L'audit réseau sera mené à l'aide de l'outil Nmap, en deux temps ; la version exacte de l'outil sera précisée lors de l'exécution. La première phase sera une découverte d'hôtes actifs par requête ARP (\texttt{nmap -sn}), sans interrogation de services, sur le segment de données du service informatique. La seconde consistera en un balayage ciblé sur sept ports associés à des services web et de vidéosurveillance (80, 443, 554, 8000, 8080, 8899, 37777), avec détection de version des services (\texttt{nmap -sV}), afin d'identifier la nature des équipements exposant un service sur ces ports. Aucun test d'intrusion, de force brute ni capture de trafic ne sera réalisé : la découverte se limitera aux services exposés, conformément au cadre éthique défini plus haut. Les sorties complètes des deux phases seront archivées sans modification en annexe~E.

\subsection{Nouvel audit programmé}

Les résultats d'un premier audit réalisé pendant le stage ne sont pas repris dans ce mémoire : un nouvel audit réseau doit être conduit ultérieurement, selon la méthode ci-dessus, pour établir un état des lieux à jour. Aucun résultat, aucune adresse IP, aucun port détecté ni aucune interprétation n'est donc présenté par anticipation : les constats, le tableau de synthèse par constructeur d'équipement et la liste des services exposés seront rédigés exclusivement à partir des sorties effectivement recueillies. \acompleter{date et périmètre exacts du nouvel audit, à fixer avec le tuteur ; tableau de synthèse des constructeurs et nombre d'hôtes, à extraire des futures sorties Nmap}

\subsection{Architecture physique du réseau}

Le service a fourni un schéma de l'architecture du réseau du CNHU-HKM\cite{archi}, dont la date exacte reste à préciser. Il montre une topologie hiérarchique : un routeur Cisco relié à Internet, un switch fédérateur (Dell) qui sert de cœur, puis des liaisons en fibre optique vers des armoires de brassage et des switches d'accès répartis par service (urgences, soins, imagerie, laboratoire, banque de sang, administration, pharmacies, etc.). Les postes et des points d'accès Wi-Fi s'y raccordent, et un serveur applicatif est relié au cœur (figure~\ref{fig:existant}). Le schéma distingue par couleur les switches Cisco existants de ceux à remplacer, par un Cisco simple ou par un Cisco à port SFP ; la majorité des switches d'accès y est marquée à remplacer.

\begin{figure}[htbp]
\centering
\small
\begin{tikzpicture}[
  box/.style={draw, rounded corners=2pt, align=center, minimum height=8mm, text width=3.4cm, font=\scriptsize},
  node distance=7mm and 6mm]
  \node[box] (net) {Internet};
  \node[box, right=of net] (rtr) {Routeur Cisco};
  \node[box, right=of rtr, fill=gray!25] (core) {Switch fédérateur (Dell)};
  \node[box, below=14mm of rtr] (srv) {Serveur applicatif};
  \node[box, below=14mm of core] (ex) {Armoires avec switches Cisco existants};
  \node[box, right=of ex] (rem) {Switches d'accès à remplacer (services)};
  \node[box, below=10mm of ex] (ap) {Points d'accès Wi-Fi};
  \node[box, dashed, below=14mm of srv] (sonde) {Sonde IDS (port miroir)\\ \tiny emplacement proposé};
  \draw (net) -- (rtr) -- (core);
  \draw (core) -- node[left, font=\tiny]{FO} (srv);
  \draw (core) -- node[left, font=\tiny]{FO} (ex);
  \draw (core.south east) -- ++(0,-0.4) -| node[above, pos=0.2, font=\tiny]{FO} (rem);
  \draw (ex) -- (ap);
  \draw[dashed] (core.south west) -- ++(-0.9,-0.5) |- (sonde.east);
\end{tikzpicture}
\caption{Architecture physique simplifiée, redessinée d'après le schéma du service\cite{archi}. L'emplacement de la sonde est une proposition}
\label{fig:existant}
\end{figure}

Quatre enseignements se dégagent pour ce travail.
\begin{itemize}
\item \textbf{Un point de convergence unique.} Le trafic entre services transite par le switch fédérateur et le routeur. C'est un emplacement d'observation idéal : un port miroir sur le fédérateur donne à une sonde passive la vue sur l'essentiel des flux, sans effet sur le trafic. C'est aussi un point de défaillance unique, ce qui pèse sur le choix d'un mode en coupure.
\item \textbf{Aucun pare-feu, VLAN ni adressage n'apparaît sur le schéma.} Cette absence de représentation ne préjuge pas de l'existence de ces éléments : le nouvel audit et les échanges avec le service devront lever cette ambiguïté.
\item \textbf{Un parc de switches hétérogène et en renouvellement.} Les capacités (port miroir, VLAN, contrôle ARP) dépendent du modèle, ce qui conditionne la faisabilité de la sonde et de la segmentation. \acompleter{modèle exact du switch fédérateur Dell et prise en charge du port miroir}
\item \textbf{Une limite de visibilité.} Le trafic entre deux postes d'un même switch d'accès ne traverse pas le cœur et échappe à une sonde placée à cet endroit.
\end{itemize}
La date exacte du schéma restant à préciser, il doit être validé par le service avant d'être exploité comme référence ; le nouvel audit permettra notamment de vérifier le nombre et le positionnement des points d'accès Wi-Fi.

\subsection{Inventaire, flux et usages}

\acompleter{inventaire détaillé par zone fonctionnelle, type d'équipement et criticité estimée, à établir après le nouvel audit en croisant ses sorties avec la documentation interne et les échanges avec le personnel technique}

\acompleter{flux métiers connus du service informatique (accès au SIH, à l'imagerie, à la messagerie), décrits à partir des échanges avec le personnel, même sans capture de trafic}

\section{Analyse de risques simplifiée}

L'analyse s'inspire d'une version simplifiée de la méthode EBIOS Risk Manager : identification des biens essentiels, des scénarios redoutés, puis cotation de la vraisemblance et de la gravité sur une échelle de 1 (faible) à 4 (critique). Le tableau~\ref{tab:risques} présente des scénarios de risque anticipés à partir des menaces documentées au chapitre~\ref{chap:1} et de l'architecture schématisée plus haut ; ils devront être confirmés ou ajustés à la lumière du nouvel audit. Les cotations sont une proposition de l'auteur, \acompleter{à valider avec le tuteur}.

\begin{table}[htbp]
\centering
\small
\begin{tabularx}{\textwidth}{lXccc}
\hline
\textbf{Réf.} & \textbf{Scénario de risque} & \textbf{Vrais.} & \textbf{Gravité} & \textbf{Niveau} \\
\hline
R1 & Propagation d'un rançongiciel depuis un poste administratif vers les serveurs applicatifs, faute de segmentation & 3 & 4 & 12 \\
R2 & Compromission d'un équipement biomédical, faute d'authentification robuste et d'isolement & 2 & 4 & 8 \\
R3 & Intrusion non détectée, faute de supervision centralisée & 3 & 3 & 9 \\
R4 & Accès d'un visiteur ou d'un équipement non maîtrisé au réseau interne & 2 & 3 & 6 \\
\hline
\end{tabularx}
\caption{Scénarios de risque (niveau = vraisemblance $\times$ gravité)}
\label{tab:risques}
\end{table}

La hiérarchisation place la propagation latérale (R1) en tête, suivie de l'absence de détection (R3). Un IDS/IPS répond directement à ces deux risques : il repère les comportements de propagation et comble l'absence de détection. Il justifie le choix du système étudié dans ce mémoire.

\section{Étude comparative des solutions IDS/IPS}

\subsection{Critères}

Les solutions ont été comparées selon des critères issus des contraintes du contexte : coût de licence nul, prise en charge des modes IDS et IPS, qualité et disponibilité des règles de détection, intégration à la plateforme de pare-feu, format de journalisation exploitable, maturité et documentation, et charge de déploiement sur un laboratoire aux ressources limitées.

\subsection{Moteurs de détection}

\begin{table}[htbp]
\centering
\small
\begin{tabularx}{\textwidth}{lXXX}
\hline
\textbf{Critère} & \textbf{Suricata} & \textbf{Snort} & \textbf{Zeek} \\
\hline
Nature & moteur IDS/IPS à règles & moteur IDS/IPS à règles & analyseur de trafic et de comportements \\
Modes IDS et IPS & oui\cite{suricatadoc} & oui & surtout analyse, pas conçu comme IPS \\
Règles & ET Open et règles locales & règles communautaires et locales & scripts d'analyse \\
Intégration à OPNsense & native\cite{opnsensedoc} & non native & non \\
Journal structuré & \texttt{eve.json} & formats propres & journaux structurés \\
\hline
\end{tabularx}
\caption{Comparaison des moteurs de détection (appréciations de l'auteur)}
\label{tab:comparatif}
\end{table}

Comme le synthétise le tableau~\ref{tab:comparatif}, Suricata est retenu : il offre les modes IDS et IPS, un large jeu de règles ouvert, un journal structuré facile à exploiter et surtout une intégration native à OPNsense, ce qui limite la charge de déploiement. Zeek, plus orienté vers l'analyse du trafic, est un complément possible plutôt qu'un substitut.

\subsection{Modes de déploiement}

Deux modes de déploiement ont été envisagés. Le premier intègre Suricata au pare-feu OPNsense, qui devient à la fois point de filtrage et point d'inspection ; il est simple à mettre en œuvre et à administrer. Le second place une sonde dédiée, alimentée par un port miroir, qui ne perturbe jamais le trafic mais ne peut pas bloquer. Le premier est retenu pour ce travail, car il permet de passer de l'IDS à l'IPS sans changer d'architecture. Le second reste adapté à une phase d'observation sur le réseau de production.

\subsection{Pare-feu et supervision}

OPNsense est retenu comme plateforme hôte pour son caractère open source, son filtrage à état et son intégration de Suricata\cite{opnsensedoc}. Wazuh est retenu comme outil complémentaire de visualisation des alertes, pour ses agents, sa corrélation et sa réponse active dans un seul produit open source\cite{wazuhdoc}.

\section{Conception du système de détection}\label{sec:conception}

\subsection{Zones surveillées et plan d'adressage}

Le système inspecte le trafic qui circule entre des zones fonctionnelles. À partir de la segmentation hospitalière décrite au chapitre~\ref{chap:1} et de l'architecture documentée, cinq zones sont retenues, chacune portée par un VLAN ; leur calage définitif sur les réseaux réels interviendra après le nouvel audit. Le plan d'adressage ci-dessous est celui du laboratoire, distinct du plan d'adressage de production, qui sera documenté à l'issue du nouvel audit. Le pare-feu en est la passerelle (adresse .1 de chaque réseau).

\begin{table}[htbp]
\centering
\small
\begin{tabular}{llll}
\hline
\textbf{VLAN} & \textbf{Zone} & \textbf{Réseau} & \textbf{Contenu type} \\
\hline
10 & Administration & 10.10.10.0/24 & postes administratifs, poste d'administration \\
20 & Soins & 10.10.20.0/24 & postes des unités de soins \\
30 & Biomédical & 10.10.30.0/24 & équipements biomédicaux, imagerie \\
40 & Serveurs & 10.10.40.0/24 & serveurs applicatifs, annuaire, serveur Wazuh \\
50 & Invités & 10.10.50.0/24 & visiteurs, accès Internet seul \\
\hline
\end{tabular}
\caption{Zones surveillées et plan d'adressage du laboratoire (proposition)}
\label{tab:vlan}
\end{table}

\subsection{Positionnement de la sonde}

Dans l'architecture cible, Suricata est placé sur le pare-feu OPNsense, par lequel transitent tous les flux entre zones et vers l'extérieur (figure~\ref{fig:cible}). Il inspecte ainsi les flux inter-zones, les plus révélateurs d'une propagation latérale, ainsi que le trafic entrant et sortant.

Sur le réseau réel, le schéma de l'architecture fait ressortir deux emplacements possibles. Le premier est un port miroir du switch fédérateur, qui alimente une sonde passive : il ne perturbe pas le trafic et convient à la phase IDS. Le second est l'insertion de la sonde en coupure entre le routeur et le switch fédérateur : il permet le blocage, mais fait de la sonde un point de défaillance et impose un mécanisme de contournement ou une redondance. Le laboratoire reproduit le second cas, car il permet de tester l'IDS comme l'IPS.

\begin{figure}[htbp]
\centering
\small
\begin{tikzpicture}[
  zone/.style={draw, rounded corners=2pt, align=center, minimum height=9mm, text width=2.2cm, font=\scriptsize},
  node distance=3mm]
  \node[zone, fill=gray!25, text width=4.2cm] (fw) {Pare-feu OPNsense\\ \tiny Suricata : sonde IDS puis IPS};
  \node[zone, above=8mm of fw, text width=3.6cm] (wan) {Réseau externe / Internet};
  \node[zone, below=14mm of fw, xshift=-5.4cm] (adm) {VLAN 10\\Administration};
  \node[zone, right=of adm] (soi) {VLAN 20\\Soins};
  \node[zone, right=of soi] (bio) {VLAN 30\\Biomédical};
  \node[zone, right=of bio] (srv) {VLAN 40\\Serveurs\\ \tiny + Wazuh};
  \node[zone, right=of srv] (inv) {VLAN 50\\Invités};
  \draw (wan) -- (fw);
  \foreach \z in {adm,soi,bio,srv,inv} \draw (fw) -- (\z);
\end{tikzpicture}
\caption{Architecture cible : Suricata inspecte les flux entre zones sur le pare-feu. À remplacer par un schéma draw.io si souhaité}
\label{fig:cible}
\end{figure}

\subsection{Flux de référence}

Pour écrire des règles pertinentes, il faut connaître les flux attendus. Le tableau~\ref{tab:flux} liste les flux légitimes entre zones ; ils servent de base aux règles de détection, car tout flux hors de cette liste est suspect. Les ports applicatifs du SIH et des équipements biomédicaux sont \acompleter{à confirmer auprès du service}. Le même tableau fonde la politique de filtrage du pare-feu, qui applique le principe « tout interdit sauf exception »\cite{nist80041}.

\begin{table}[htbp]
\centering
\small
\begin{tabularx}{\textwidth}{llXl}
\hline
\textbf{Source} & \textbf{Destination} & \textbf{Service} & \textbf{Décision} \\
\hline
Administration & Serveurs & HTTPS 443, applications du SIH & Autoriser \\
Soins & Serveurs & applications du SIH, DICOM 104, HL7 (MLLP) & Autoriser \\
Biomédical & Serveurs (imagerie) & DICOM uniquement & Autoriser \\
Biomédical & Autres zones & tout & Interdire \\
Administration & Toutes zones & SSH 22, RDP 3389 depuis les postes d'administration & Autoriser, journaliser \\
Toutes zones & Serveurs (annuaire) & DNS 53, DHCP 67/68 & Autoriser \\
Hôtes supervisés & Serveur Wazuh & agent 1514, enregistrement 1515 & Autoriser \\
Pare-feu & Serveur Wazuh & syslog 514 & Autoriser \\
Invités & Internet & DNS, HTTP, HTTPS & Autoriser \\
Invités & Zones internes & tout & Interdire \\
\hline
\end{tabularx}
\caption{Flux de référence entre zones (proposition)}
\label{tab:flux}
\end{table}

\subsection{Stratégie de déploiement et gestion des règles}

La mise en service est progressive. Suricata fonctionne d'abord en mode IDS, ce qui permet d'observer les alertes sans risque pour le trafic. Après réglage des faux positifs, certaines règles à haute confiance passent en mode IPS (action \texttt{drop}). Le jeu de règles combine ET Open et des règles locales couvrant les flux DICOM et HL7, les accès d'administration anormaux et les comportements de balayage ou de force brute. Chaque alerte est journalisée dans \texttt{eve.json}, conservée pour analyse et relayée vers Wazuh.

\subsection{Éléments complémentaires}

La segmentation par zones et la supervision par Wazuh complètent le système sans en être l'objet : la première donne à la sonde de la visibilité sur les flux inter-zones, la seconde facilite la lecture des alertes. Elles sont mises en place dans le laboratoire à ce titre, et leur généralisation au réseau de production relève des recommandations.

\addcontentsline{toc}{section}{Conclusion}
\section*{Conclusion}

La problématique reste fondée sur les besoins exprimés par le service et sur les menaces hospitalières documentées au chapitre~\ref{chap:1} ; le nouvel audit prévu permettra de la confirmer par des mesures. En l'état, la proposition d'architecture repose sur Suricata, placé sur OPNsense, inspectant les flux entre cinq zones, d'abord en IDS puis en IPS. Le chapitre suivant décrit le protocole de déploiement en laboratoire et la méthode d'évaluation prévue, dont l'exécution reste à conduire.
